\begin{abstract}
Dense retrieval often places the query encoder on the online critical path while document encoding is amortized offline. This asymmetry motivates query-side compression, but it is unclear how much retrieval quality survives when only the query tower is aggressively reduced. We study this question with a frozen zero-shot \bgebase\ teacher (\texttt{BAAI/bge-base-en-v1.5}), asymmetric 4-layer and 2-layer students, and a controlled MS MARCO-derived benchmark with explicit latency measurement. Naive truncation causes severe retrieval collapse. Among low-cost interventions, query-side mean pooling is the first clear positive change, whereas a query-side projection head hurts and lightweight distillation is nearly neutral. A later systems comparison shows that IVF, but not HNSW, provides a plausible search-speed tradeoff relative to exact flat search. The resulting picture is therefore that, shallow asymmetric query compression is brittle under simple truncation, and internal design choices matter before strong efficiency claims can be made. You may find the code at: \url{https://github.com/AngadSingh22/FastQueryDR}.
\end{abstract}
